% !PRINTED 438
\weberpage{438}
From these formulae one easily obtains, by (5), (8), (14):
\weq{18}{\begin{aligned}
\omega\Theta &= \begin{pmatrix} \gamma, & \alpha\varepsilon^{-2}, & \beta\varepsilon \\ \alpha\varepsilon^2, & \beta, & \gamma\varepsilon^3 \\ \beta\varepsilon^{-1}, & \gamma\varepsilon^{-3}, & \alpha \end{pmatrix}, \quad
\omega\Theta^2 = \begin{pmatrix} \gamma, & \alpha\varepsilon, & \beta\varepsilon^3 \\ \alpha\varepsilon^{-1}, & \beta, & \gamma\varepsilon^2 \\ \beta\varepsilon^{-3}, & \gamma\varepsilon^{-2}, & \alpha \end{pmatrix}, \\
\Theta\omega &= \omega\Theta^3 = \begin{pmatrix} \alpha, & \beta\varepsilon^{-1}, & \gamma\varepsilon^{-3} \\ \beta\varepsilon, & \gamma, & \alpha\varepsilon^{-2} \\ \gamma\varepsilon^3, & \alpha\varepsilon^2, & \beta \end{pmatrix}.
\end{aligned}}

After this it is now very simple to confirm the characteristic relations of Theorem I., \S.\,72 for our group by actual computation, and with this the ternary group of degree 168 is therefore established\footnote{Cf. F. \gesperrt{Klein}, ``Ueber die Transformation siebenter Ordnung der elliptischen Functionen''. ``Ueber die Auflösung gewisser Gleichungen vom 7\textsuperscript{ten} und 8\textsuperscript{ten} Grade''. Mathem. Annalen, Bd. 14 u. 15. \gesperrt{Klein-Fricke}, Vorlesungen über Modul-Functionen, Bd. I, dritter Abschnitt, Capitel VII. \gesperrt{Gordan}, ``Ueber Gleichungen 7\textsuperscript{ten} Grades mit einer Gruppe von 168 Substitutionen''. Mathem. Annalen, Bd. 20, 25.}).

\begin{center}
\S.\,116.\\[2pt]
\gesperrt{Poles and Axes of the Ternary Groups.}
\end{center}

We shall now, for the sake of brevity, make use of a geometrical mode of expression, by regarding the variables $x_1, x_2, x_3$ as triangular coordinates of a point $x$ in a plane, although imaginary values of $x_1, x_2, x_3$ are also to be admissible.

If $A$ denotes a ternary linear substitution with determinant 1, and if
\weq{1}{(y) = A(x),}
then $y_1, y_2, y_3$ are the coordinates of a second point $(y)$, referred to the same coordinate system, which is derived from $(x)$ by the substitution $A$.

When $x$ traverses a straight line, then $y$ also traverses a straight line, and if the equations of these two lines are
\weq{2}{u_1x_1 + u_2x_2 + u_3x_3 = 0, \quad v_1y_1 + v_2y_2 + v_3y_3 = 0}
then it follows, by substituting (1) in (2), that $u_1, u_2, u_3$ are connected with $v_1, v_2, v_3$ by the \gesperrt{transposed substitution} of $A$:
\[ (u) = A_1(v) \]
%CONT>>
